% !PRINTED 324
%<<CONT
Stelle von $\varphi$ gesetzt werden. Es bleiben die beiden willkürlichen
Grössen $c$ und $d$ darin.

Setzen wir der eleganteren Bezeichnung wegen noch $c = -\eta$,
$d = \xi$, so erhalten wir also:

\weq{8}{S\left(\frac{\xi\beta - \eta\alpha}{x\beta - y\alpha}\right)^2 = \frac{[\xi f'(x) + \eta f'(y)]^2 - (x\eta - \xi y)^2 H(x,y)}{nf(x,y)^2},}

und in dieser Form bleibt sowohl die rechte als die linke Seite
vollkommen ungeändert, wenn wir auf die Variablenpaare $x, y$;
$\xi, \eta$; $\alpha, \beta$ gleichzeitig eine beliebige lineare Transformation
anwenden.

Die Formel (8) ist eine identische; sie gilt für jede Function
$f(x, y)$ und für alle Werthe der Variablen $x, y$; $\xi, \eta$. Jetzt aber
machen wir die Voraussetzung, dass $f(x, y)$ nur reelle Wurzeln habe,
dass also die $\alpha, \beta$ reell seien, und wir setzen auch $\xi, \eta$ und
$x, y$ als reell voraus. Beide Seiten der Gleichung (8) sind dann
wesentlich positiv. Wir bezeichnen ihren gemeinsamen Werth durch
$P$, und bestimmen nun eine Grösse $X:Y$ durch die quadratische
Gleichung
\[
P = \left(\frac{X\eta - Y\xi}{Xy - Yx}\right)^2,
\]
oder

\weq{9}{\Phi = (Xy - Yx)^2 P - (X\eta - Y\xi)^2 = 0.}

Diese Gleichung hat zwei reelle Wurzeln, die man aus
\[
X\eta - Y\xi = \pm \sqrt{P}\,(Xy - Yx)
\]
erhält. Ueber die Lage der Wurzeln dieser Gleichung können wir aber
Folgendes aussagen:

Setzen wir $X = x$, $Y = y$, so wird $\Phi$ negativ.

Setzen wir aber $X = \alpha$, $Y = \beta$, wenn $\alpha:\beta$ irgend eine
der Wurzeln von $f = 0$ ist, so ist
\[
\left(\frac{X\eta - Y\xi}{Xy - Yx}\right)^2 = \left(\frac{\xi\beta - \eta\alpha}{x\beta - y\alpha}\right)^2,
\]
also gleich \gesperrt{einem Gliede} der Summe
\[
P = S\left(\frac{\xi\beta - \eta\alpha}{x\beta - y\alpha}\right)^2
\]
und daher kleiner als $P$. Daraus folgt, dass $\Phi$ für $X:Y = \alpha:\beta$
\gesperrt{positiv ist}.

Wenn wir nun die Werthe der sämmtlichen $\alpha:\beta$ und $x:y$, der
Grösse nach cyklisch geordnet, etwa auf einem Kreise anordnen,
%CONT>>
